\documentclass[10pt,a4paper]{amsart}
\usepackage[margin=19mm]{geometry}
\usepackage{amsmath,amssymb,mathtools}
\usepackage[scr=rsfs,cal=euler]{mathalfa}
\usepackage{xfrac}
\usepackage[unicode,hidelinks,bookmarks=false]{hyperref}
\newcommand{\ie}{i.e.\ }
\newcommand{\cf}{cf.\ }
\newcommand{\resp}{respectively\ }
\newcommand{\Subsection}{\subsection}
\providecommand{\nobreakdash}{\nobreak\hbox{-}}
\setlength{\emergencystretch}{2em}
\multlinegap0pt
\usepackage{bm}
\usepackage{bbm}
\usepackage{dsfont}
\usepackage{xspace}
\usepackage{cancel}
\usepackage{mathtools}
\usepackage{amsthm}
\makeatletter
\@ifundefined{theorem}{\newtheorem{theorem}{Theorem}}{}
\@ifundefined{corollary}{\newtheorem{corollary}[theorem]{Corollary}}{}
\@ifundefined{lemma}{\newtheorem{lemma}[theorem]{Lemma}}{}
\makeatother
\newcommand {\N}{\mathbb N}
\newcommand{\F}{\mathbb F}
\newcommand{\pdt}{\underline}
\title[Radicals of Biduals of Beurling Algebras Can Be Different fo]{Radicals of Biduals of Beurling Algebras Can Be Different for the Two Arens Products}
\author{Jared T. White}
\date{Source: arXiv:2602.02764v1 (CC BY 4.0)}
\begin{document}
\maketitle

\noindent\emph{Source and licence.} Radicals of Biduals of Beurling Algebras Can Be Different for the Two Arens Products. Version arXiv:2602.02764v1 (CC BY 4.0), distributed under the Creative Commons Attribution 4.0 International licence, as recorded in the arXiv licence metadata for this exact version and shown on the abstract page https://arxiv.org/abs/2602.02764v1. Full rights notice: licenses/arxiv-2602.02764-CC-BY-4.0.txt.\par\smallskip

\noindent\textit{Editorial scope.} Counted unit: the Lemma whose source label is \texttt{3.7} in the article (source0.tex lines 429--439, source label \texttt{3.7}), delivered together with the article's own complete original proof of it (source0.tex lines 442--560, one \texttt{proof} environment of 119 lines). No mathematics is written, repaired or re-derived: statement and proof are the article's own text line for line. Required context retained uncounted: 3.0 (lines 299--301); 3.1a (lines 327--330); 3.6 (lines 400--405). Every definition, display and label of those lines is retained unchanged. Omitted from this package and not redistributed: 1--298, 302--326, 331--399, 406--428, 440--441, 561--1028. Nothing inside a retained excerpt is omitted or altered: every retained body is delivered exactly as the article's lines are written, so no omission receipt of any kind accompanies this package. Citations: the retained excerpts carry no citation command at all, so no bibliography entry of the article is transcribed and no reference is redistributed. Reference adaptation: none was needed and none was made; every reference inside the delivered unit resolves inside it, so no unresolved reference or citation is produced. Printed slips kept literal: no printed word, symbol, hypothesis or number of the counted statement or of its proof was altered. Layout and class: the article's own class is \texttt{amsart} with packages including a4wide, latexsym, mathrsfs, fontenc, enumitem, amssymb, amsmath, fancyhdr, bbm, setspace, mathabx; the wrapper is the lane's pinned \texttt{amsart} editorial wrapper at 10pt on A4, which restates the theorem-like environments the retained text needs and adds the article's own macro definitions that the retained text needs (\texttt{\textbackslash F}, \texttt{\textbackslash N}, \texttt{\textbackslash pdt}), with extra wrapper packages (bm, bbm, dsfont, xspace, cancel, mathtools). The delivered numbering is the unit's own. Assets: no retained span contains a figure environment, a graphics-inclusion command, a PDF-page inclusion, a table or a TikZ picture; no image file is copied into this package, the package declares no asset dependency and no image file is redistributed with it. Licence: version arXiv:2602.02764v1 of this article is distributed under the Creative Commons Attribution 4.0 International licence, as recorded in the arXiv licence metadata for this exact version and shown on the abstract page https://arxiv.org/abs/2602.02764v1. A full rights notice is delivered at licenses/arxiv-2602.02764-CC-BY-4.0.txt. Characters: every retained excerpt is pure ASCII, so the delivered engine, XeLaTeX with its default Latin Modern text font, reports no missing character.
	\begin{lemma} 	\label{3.0}
		Let $u \in \F_3$. Then $|u| \geq |\varphi_a(u)| + |\varphi_b(u)| + |\varphi_c(u)|.$
	\end{lemma}

	\begin{corollary}		\label{3.1a}
		Let $n \in \N$ and let $v \in \F_3$ with $|v| \leq 5^n$. As a reduced word $x(v,n)$ has the form
		$a^k w b^{5^{2n-1}} w^{-1} a^{5^{2n}-k}b^{5^{2n}}$, and  $w \in \F_3$ and $k \in \N$ satisfy $|w|+k \leq 5^n$. 
	\end{corollary}

	\begin{lemma}		\label{3.6}
		Let $j \in \N$, and let $v \in \F_3$ satisfy $|v| \leq 5^j$. Write 
		$x(v,j) = s_1s_2 \cdots s_m$ as a reduced word, where $m = |x(v,j)|$ and $s_1, \ldots, s_m \in S$.
		Suppose that $z$ is an initial block of these letters, and that $y$ is the rest (so that in particular $x(v,j) = zy$).
		Suppose further that $|z| < \frac{1}{2} |x(v,j)|$. Then $\theta(y) \geq |z|$. 
	\end{lemma}

	\begin{lemma}		\label{3.7} 
		Let $n \in \N$ and let $x_1, \ldots, x_n \in X_\infty$. Suppose that the expression
		\begin{equation}	\label{eq3.15}
			\pdt{d}_1x_1 \pdt{d}_2 \cdots \pdt{d}_n x_n
		\end{equation}
		is super $X$-minimal, where $x_i \in X_\infty$ and $d_{ij} \in S$ for all $i,j$. Write 
		$$x_i = a^{\nu_i}w_ib^{5^{2k_i-1}}w_i^{-1}a^{5^{2k_i} - \nu_i}b^{5^{2k_i}} \ \text{(reduced)} \quad (i = 1, \ldots, n),$$
		as in Corollary \ref{3.1a}, where $\nu_1, \ldots, \nu_n \in \N_0$, $w_1, \ldots, w_n \in \F_3$, and $k_1, \ldots, k_n \in \N$.			
		Fix a sequence of cancellations that reduces the expression \eqref{eq3.15}.
		Then if any terms from $a^{5^{2k_n} - \nu_n}$ are cancelled, then they are cancelled by letters from $\pdt{d}_n$.	
	\end{lemma}

	\begin{proof}	
		Suppose that $r$ of the letters from $a^{5^{2k_n} - \nu_n}$ are cancelled. Then there must be some contiguous block of letters at the end of the string
		\begin{equation}		\label{eq3.14}
			\pdt{d}_1x_1 \pdt{d}_2 \cdots x_{n-1} \pdt{d}_n \, a^{\nu_n} w_n b^{5^{2k_n-1}} w_n^{-1}
		\end{equation}
		that reduces to $a^{-r}$. Our proof shall proceed by dividing into cases according to the various possibilities for where this block can lie. Only Case 1 and Subcase 4C at the very end do not result in a contradiction.
		
		\vskip 2mm
		\noindent
		\underline{Case 1.} The block comes entirely from $\pdt{d}_n a^{\nu_n} w_n b^{5^{2k_n-1}} w_n^{-1}$.
		
		In this case, we see that $\pdt{d}_n$ must cancel $b^{5^{2k_n-1}}$, and hence (as \eqref{eq3.15} is super $X$-minimal) $w_n=1$ and $\nu_n = 0$. It is then clear that $\pdt{d}_n = y a^{-r}b^{-5^{2k_n-1}}$ (reduced), for some $y \in \F_3$, and  in particular that the cancellations all come from the $a^{-1}$'s in $\pdt{d}_n$.
		
		\vskip 2mm
		\noindent
		\underline{Case 2.} The block ends somewhere within $\pdt{d}_m$, for some $m<n$. 
		
		%Say the block terminates at $d_{m, l}$ for some index $l$, and write $\pdt{d}_m'' = d_{m,1}d_{m,2} \cdots d_{m, l-1}$, and $\pdt{d}_m' = d_{m,l}d_{m, l+1} \cdots d_{m,p_m}$. (I haven't defined $p_m$.)
		
		%In this case write $\pdt{d}_m = \pdt{d}_m''\pdt{d}_m'$ (reduced), where $\pdt{d}_m'$ is the part of $\pdt{d}_m$ that belongs to the block.
		
		In this case, divide the sequence $d_m$ into into an initial part $d_m'$, and the rest $d_m''$, so that $\pdt{d}_m''$ is the part of $\pdt{d}_m$ that belongs to the block. We have
		$$ a^{-r} = \pdt{d}_m''x_m \pdt{d}_{m+1} \cdots x_{n-1} \pdt{d}_n \, a^{\nu_n} w_n b^{5^{2k_n-1}} w_n^{-1}.$$
		Taking $\varphi_a$ of both sides, and noting that $\varphi_a(x_i) > 1$ for all $i$, we have
		\begin{align*}
			-r &= \varphi_a(\pdt{d}_m'')  +\sum_{i = m+1}^n \varphi_a(\pdt{d}_i) 
			+\sum_{i=m}^{n-1} \varphi_a(x_i) + \nu_n\\
			&> \varphi_a(\pdt{d}_m'')  +\sum_{i = m+1}^n \varphi_a(\pdt{d}_i) + (n-m) +\nu_n,
		\end{align*}
		which implies that
		$$ -\varphi_a(\pdt{d}_m'')  - \sum_{i = m+1}^n \varphi_a(\pdt{d}_i) > r +(n-m) +\nu_n.$$
		Similarly
		$$ -\varphi_b(\pdt{d}_m'')  -\sum_{i = m+1}^n \varphi_b(\pdt{d}_i) > (n-m) + 5^{2k_n-1}.$$
		By Lemma \ref{3.0} these last two equations together imply that  
		\begin{align*}
			|\pdt{d}_m''| + \sum_{i=m+1}^n |\pdt{d}_i | &\geq
			|\varphi_a(\pdt{d}_m'')| + |\varphi_b(\pdt{d}_m'')| + \sum_{i=m+1}^n |\varphi_a(\pdt{d}_i)| +  \sum_{i=m+1}^n |\varphi_b(\pdt{d}_i)|\\
			&\geq r+ \nu_n  + 5^{2k_n-1} + 2(n-m)> r + \nu_n +  5^{2k_n-1} + 1.
		\end{align*}
		The $X$-word
		$$\pdt{d}_m''x_m \pdt{d}_{m+1} \cdots \pdt{d}_n x_n$$
		is equal to $a^{5^{2k_n}-r - \nu_n}b^{5^{2k_n}}$ as a reduced word, and contains 
		$$(n-m) + |\pdt{d}_m''| + \sum_{i=m+1}^n |\pdt{d}_i | \geq 2 + |\pdt{d}_m''| + \sum_{i=m+1}^n |\pdt{d}_i | > r + \nu_n + 5^{2k_n-1} + 3$$
		many symbols from $X$,
		and hence replacing this string in \eqref{eq3.15} by
		$$(a^{-1})^{r +\nu_n} (b^{-1})^{5^{2k_n-1}} x(1,k_n)$$
		would allow us to represent the same element of $\F_3$ whilst using fewer elements of $X$, a contradiction.
		
		\vskip 2mm
		\noindent
		\underline{Case 3.} The block ends either between $x_{m-1}$ and $\pdt{d}_m$, or between $\pdt{d}_m$ and $x_m$, for some $m<n$.
		
		In this case, repeat the argument of Case 2., with $\pdt{d}_m'' = \pdt{d}_m$ and $\pdt{d}_m' = 1$, or $\pdt{d}_m'' = 1$ and $\pdt{d}_m' = \pdt{d}_m$, respectively.
		
		\vskip 2mm
		\noindent
		\underline{Case 4.} The block ends part way through $x_m$, for some $m<n$.
		
		Write 
		$$x_m = z_1z_2 \cdots z_R \, y_1 y_2 \cdots y_Q \ \text{(reduced)},$$
		for some $Q, R \in \N$ such that $y_1 y_2 \cdots y_Q$ is part of the block, and $z_1 z_2 \cdots z_R$ isn't. 
		The expression $x_m \pdt{d}_{m+1} \cdots \pdt{d}_n x_n$ has $n-m + 1 + \sum_{i=m+1}^n |\pdt{d}_i|$ symbols from $X$ and as an element of $\F_3$ is equal to $\pdt{z} a^{5^{2k_n}- r- \nu_n}b^{5^{2k_n}}$.
		
		We have
		$$\pdt{y} \, \pdt{d}_{m+1} x_{m+1} \cdots \pdt{d}_n  \, w_n b^{5^{2k_n-1}} w_n^{-1} = a^{-r - \nu_n}.$$
		Applying $\varphi_a$ to both sides gives
		$$\varphi_a(\pdt{y})  + \sum_{i=m+1}^n \varphi_a(\pdt{d}_i) + \sum_{i=m+1}^{n-1} \varphi_a(x_i)  = -r - \nu_n,$$
		and hence
		\begin{equation}		\label{eq3.16}
			\sum_{i=m+1}^n \varphi_a(\pdt{d}_i) = -r -\nu_n -\varphi_a(\pdt{y}) -\sum_{i=m+1}^{n-1} \varphi_a(x_i) \leq -r -\nu_n - \varphi_a(\pdt{y}). 
		\end{equation} 
		Similarly 
		\begin{equation}		\label{eq3.17}
			\sum_{i=m+1}^n \varphi_b(\pdt{d}_i) \leq -5^{2k_n-1} - \varphi_b(\pdt{y}).
		\end{equation}
		Using Lemma \ref{3.0} and noting that $\varphi_a(\pdt{y}), \varphi_b(\pdt{y}) \geq 0$, Equations \eqref{eq3.16} and \eqref{eq3.17} together imply that 
		\begin{equation}		\label{eq3.13}
			\sum_{i=m+1}^n |\pdt{d}_i| \geq \varphi_a(\pdt{y}) + \varphi_b(\pdt{y})  +r +\nu_n + 5^{2k_n-1} = \theta(\pdt{y})  +r +\nu_n + 5^{2k_n-1}.
		\end{equation}
		We now divide into subcases depending on whether or not $|\pdt{y}|\leq \frac{1}{2}|x_m|$.
		
		
		\vskip 2mm
		\noindent
		\underline{Subcase 4A.} $|\pdt{y}| > \frac{1}{2}|x_m|$.
		
		In this case Lemma \ref{3.6} tells us that $\theta(\pdt{y}) \geq |\pdt{z}|$, and so \eqref{eq3.13} implies that 
		$$\sum_{i=m+1}^n |\pdt{d}_i| \geq |\pdt{z}| + 5^{2k_n-1} + r +\nu_n = R  + r +\nu_n + 5^{2k_n-1} ,$$
		so that the expression $x_m \pdt{d}_{m+1} \cdots \pdt{d}_n x_n$ has at least 
		$$2 + \sum_{i=m+1}^n |\pdt{d}_i| \geq 2+ R + r +\nu_n + 5^{2k_n-1}$$
		symbols from $X$. 	
		On the other hand, the expression
		$$z_1 z_2 \cdots z_R \, (a^{-1})^{r+\nu_n} (b^{-1})^{5^{2k_n-1}} x(1, k_n)$$
		also equals $\pdt{z} a^{5^{2k_n}-r - \nu_n}b^{5^{2k_n}}$, but has only $1+ R + r +\nu_n + 5^{2k_n-1}$ symbols from $X$, a contradiction.
			
		\vskip 2mm
		\noindent
		\underline{Subcase 4B.} $|\pdt{y}| \leq \frac{1}{2}|x_m|$  and $m<n-1$.
		
		This means that $\pdt{y} | a^{5^{2k_m} - \nu_m} b^{5^{2k_m}},$ so that $\theta(\pdt{y}) = |\pdt{y}|$. As such, by \eqref{eq3.13},
		$$ \sum_{i=m+1}^n |\pdt{d}_i| \geq |\pdt{y}| + r + \nu_n + 5^{2k_n-1} = Q + r + \nu_n + 5^{2k_n-1}  ,$$
		and it follows that the expression $x_m \pdt{d}_{m-1} \cdots \pdt{d}_n x_n$ contains at least 
		$$n-m + 1 + \sum_{i = m-1}^n|\pdt{d}_i| \geq  3 + Q + r + \nu_n + 5^{2k_n-1}$$
		symbols from $X$.
		Therefore, the expression
		$$x_m y_Q^{-1} \cdots y_2^{-1} y_1^{-1}(a^{-1})^{r + \nu_n} (b^{-1})^{5^{2k_n-1}} x(1, k_n)$$
		has fewer symbols from $X$ than $x_m \pdt{d}_{m+1} \cdots \pdt{d}_n x_n$, whilst representing the same group element. This is a contradiction. 
		
		\vskip 2mm
		\noindent
		\underline{Subcase 4C.} $|\pdt{y}| \leq \frac{1}{2}|x_m|$  and $m=n-1$.
		
		This means that our block is $\pdt{y}\pdt{d}_n a^{\nu_n} w_n b^{5^{2k_n-1}} w_n^{-1}$, where again  $\pdt{y} | a^{5^{2k_m} - \nu_m} b^{5^{2k_m}}$. 
		We have 
		$$\pdt{y} \pdt{d}_n a^{\nu_n} w_n b^{5^{2k_n-1}} w_n^{-1} = a^{-r}$$
		and for this product to be equal to $a^{-r}$ something must cancel $\pdt{y}$. If terms from $w$ cancel $\pdt{y}$ then they must cancel $\pdt{d}_n$ first, meaning that there is nothing available to cancel $b^{5^{2k_n-1}}$. 
		Hence the cancellation must come from $\pdt{d}_n$ and we can write $\pdt{d}_n = \pdt{d}_n' \pdt{d}_n''$ (reduced), 
		where $\pdt{d}_n' = y_Q^{-1}\cdots y_2^{-1} y_1^{-1}$. Then  $\pdt{d}_n'' a^{\nu_n} w_n b^{5^{2k_n-1}} w_n^{-1} = a^{-r}$. Now we may argue as in Case 1 to get the desired conclusion.
	\end{proof}

\end{document}
